\documentclass[11pt,a4paper]{article}

\usepackage[T1]{fontenc}
\usepackage[utf8]{inputenc}
\usepackage{lmodern}
\usepackage[margin=0.82in]{geometry}
\usepackage{microtype}
\usepackage{amsmath,amssymb}
\usepackage{enumitem}
\usepackage{xcolor}
\usepackage{hyperref}
\usepackage{fancyhdr}

\hypersetup{colorlinks=true,linkcolor=blue!50!black,urlcolor=blue!50!black}
\setlist{nosep,leftmargin=*}
\setlength{\parindent}{0pt}
\setlength{\parskip}{0.45em}
\pagestyle{fancy}
\fancyhf{}
\lhead{LM-IDNet}
\rhead{Crucial Prototype Roadmap}
\cfoot{\thepage}
\sloppy

\newcommand{\status}{\(\square\) Not started \quad \(\square\) In progress \quad
  \(\square\) Blocked \quad \(\square\) Done}
\newcommand{\task}[4]{%
  \subsection{#1: #2}
  \textbf{Implementation.} #3

  \textbf{Verification and acceptance gate.} #4

  \textbf{Status.} \status
}
\newcommand{\artifact}[1]{\texttt{#1}}
\newcommand{\gate}[1]{\textcolor{blue!45!black}{\textbf{Phase exit gate:} #1}}

\title{LM-IDNet: Crucial Prototype Implementation Roadmap\\
\large Minimum Ordered Path to a Working IoT Anomaly Detector}
\author{Master's Thesis Engineering Plan}
\date{\today}

\begin{document}
\maketitle

\begin{abstract}
This document is the minimum ordered implementation path extracted from the
Atomic Implementation Roadmap. It ends with a working prototype that can
deterministically ingest data, produce validated time windows, fit a
Dirichlet--Multinomial model, evaluate likelihoods using the LM backend,
calibrate a threshold, and emit anomaly events. The original task identifiers
are retained for traceability. Tasks omitted from this document are additional
validation, comparative research, adaptation, forecasting, deployment, or
polishing work; they remain important for the final thesis artifact but do not
block the first working prototype.
\end{abstract}

\tableofcontents
\newpage

\section{How to Use This Roadmap}

Implement tasks strictly in document order. A task is Done only when its stated
verification passes and no previously completed verification regresses. Keep
generated evidence under \artifact{reports/} or \artifact{artifacts/}. Do not
begin the next phase until its exit gate passes.

\section{Stage 1: Executable Project Foundation}

\task{F0.02}{Adopt a source-layout package}{Define one import strategy for
\artifact{src/}; add package metadata in \artifact{pyproject.toml}; expose a
console entry point such as \artifact{lm-idnet}; remove runtime
\artifact{sys.path} mutation after migration.}{Install the project into a fresh
virtual environment in editable mode, execute \artifact{python -c "import ..."},
invoke \artifact{lm-idnet --help}, and run tests without setting
\artifact{PYTHONPATH}.}

\task{F0.04}{Centralize typed configuration}{Define a Pydantic configuration
model for paths, dates, categories, seeds, estimator settings, calibration, and
output locations. Reject unknown fields and invalid cross-field combinations.}
{Unit-test missing fields, extra fields, invalid dates, duplicate categories,
non-positive window sizes, invalid quantiles, and a known-valid configuration.
Snapshot the normalized configuration.}

\task{F0.09}{Define the testing architecture}{Configure pytest markers for
\artifact{unit}, \artifact{integration}, \artifact{numerical}, and
\artifact{slow}. Establish fixture factories for count matrices, windows,
configurations, and temporary artifact stores.}{Run each marker independently;
ensure no test is silently uncollected and that the default fast suite excludes
only explicitly marked slow or data-heavy tests.}

\task{F0.17}{Define exception and failure semantics}{Create domain exceptions
for configuration, ingestion, validation, numerical precision, convergence,
artifact compatibility, and policy rejection. Ensure commands fail closed when
model integrity is uncertain.}{Unit-test every exception mapping to CLI exit
code and message. Corrupt a model checksum and verify scoring stops rather than
continuing with a warning.}

\task{F0.18}{Version schemas and artifacts}{Assign semantic schema versions to
processed datasets, models, thresholds, events, and experiment manifests.
Reject unknown major versions.}{Round-trip each schema through JSON and reject
an unsupported future major version.}

\task{F0.06}{Standardize the command-line interface}{Implement the prototype
subcommands \artifact{preprocess}, \artifact{train}, \artifact{calibrate}, and
\artifact{score}. Make errors non-zero and machine-readable where appropriate.}
{For every implemented subcommand, test \artifact{--help}, missing arguments,
invalid configuration, and a minimal success case. Assert exit codes and
essential output fields.}

\gate{The package installs cleanly; typed configuration, tests, versioned
schemas, failure semantics, and the four prototype commands work from a fresh
environment.}

\section{Stage 2: Reproducible Data Setup}

\task{P0.01}{Inventory source captures}{Enumerate expected and present captures
with date, byte size, packet count where readable, checksum, and parser status.
Record missing or duplicate files without silently skipping them.}{Compare
inventory entries with the filesystem; assert unique checksums unless
duplication is explained; fail when a configured required capture is missing.}

\task{P0.02}{Verify capture integrity}{Use a PCAP validation tool or complete
streaming parse to detect truncation, invalid timestamps, and corrupt records.
Record recoverable warnings separately from fatal corruption.}{Parse every
configured source to end-of-file. Store counts and warnings; rerun and confirm
identical counts. Quarantine a deliberately truncated fixture and verify the
failure classification.}

\task{P0.05}{Freeze temporal partitions}{Create disjoint fit, calibration,
development-test, and final-test date lists. Preserve chronology and prohibit
random window mixing across partitions.}{Automated tests assert non-overlap,
chronological order, capture-level grouping, and that final-test identifiers are
unavailable to training and threshold code.}

\task{P0.06}{Create dataset fingerprints}{Hash the ordered file list,
individual inputs, feature taxonomy, window policy, and partition definition
into a dataset version identifier.}{Changing any file, category order, window
length, or date assignment must change the fingerprint; repeating unchanged
inputs must reproduce it.}

\task{P0.07}{Centralize deterministic seeds}{Define named seeds for simulation
and model fitting. Propagate generators explicitly instead of using global
state.}{Run the smoke experiment twice in separate processes and compare model
parameters, injected anomalies, metrics, and normalized reports exactly or
within a declared numerical tolerance.}

\task{P0.10}{Build a tiny public smoke fixture}{Create a payload-free synthetic
packet/count fixture small enough for tests, with known timestamps, categories,
silent windows, and expected outputs.}{Verify the fixture contains no real
identifiers, its license permits distribution, and its expected preprocessing
result is asserted exactly in integration tests.}

\gate{Inputs are inventoried and validated, chronological partitions and the
dataset fingerprint are frozen, and a safe deterministic smoke fixture exists.}

\section{Stage 3: Canonical Windowed Data}

\task{P1.01}{Define canonical category semantics}{Choose lowercase canonical
names and explicitly decide whether SSDP is exclusive of UDP or counted in
both. Record category order and classification priority.}{Use packets
representing TCP, ordinary UDP, SSDP/UDP:1900, ARP, and unknown protocols;
assert exactly the required counts and fixed column order.}

\task{P1.03}{Harden timestamp parsing}{Convert all packet times to
timezone-aware UTC, reject non-finite values, and preserve sufficient precision.
Record original capture time-zone assumptions.}{Test epoch boundaries,
fractional seconds, naive values, out-of-order packets, and daylight-saving
examples; assert normalized UTC values.}

\task{P1.04}{Specify window boundaries}{Implement half-open windows
\([t,t+\Delta)\), with an explicitly chosen origin and deterministic boundary
behavior.}{Place packets immediately before, on, and immediately after a
boundary and assert assignment to exactly one correct window for every
supported window length.}

\task{P1.05}{Preserve window timestamps}{Define each window record with device
ID, start/end UTC timestamps, ordered counts, total count, missingness state,
and metadata.}{Schema round-trip preserves timestamps and category order;
assert \(N_t=\sum_{k=1}^{K}x_{tk}\) for every record.}

\task{P1.06}{Represent silent windows}{Materialize elapsed capture windows with
zero total count and label them \artifact{observed-silent}; do not confuse them
with absent capture coverage.}{Use a fixture with a deliberate quiet interval
and verify all expected zero windows exist, remain ordered, and survive
JSON/matrix export.}

\task{P1.07}{Represent missing coverage}{Define a separate
\artifact{missing} state for gaps where observation is unavailable; prohibit
automatic conversion of missing coverage to zeros.}{Use a fixture with declared
capture discontinuity and assert missing windows are excluded or handled
according to policy, never counted as observed silence.}

\task{P1.08}{Validate count values}{Reject booleans, negatives, fractional
values, NaN, infinity, and integer overflow; use a documented integer width.}
{Parameterized tests exercise every invalid type and maximum supported values.
Valid arrays retain exact integer values through all exports.}

\task{P1.09}{Validate temporal streams}{Check monotonicity, duplicate windows,
overlap, expected duration, device continuity, and category-schema consistency.}
{Inject one defect of each type and assert a specific validation error containing
dataset and window identifiers. A valid multi-day stream passes.}

\task{P1.10}{Make PCAP parsing resource-safe}{Stream packets with bounded
memory, close readers on success and failure, constrain pathological packet
handling, and count unsupported packets.}{Measure memory over a representative
capture, inject parser exceptions, and assert file descriptors close and
partial outputs are not promoted.}

\task{P1.12}{Export canonical JSON}{Define canonical JSON with schema version,
dataset fingerprint, metadata, window records, and checksums.}{Round-trip a
multi-day fixture and compare every field. Verify stable serialization and
checksum across repeated exports.}

\task{P1.13}{Export modeling matrices}{Export timestamps, missingness mask,
category columns, totals, device ID, and partition without losing provenance.}
{Load canonical JSON and matrix output and assert row count, order, counts,
timestamps, and metadata agree exactly.}

\task{P1.14}{Load matrices by partition}{Implement a loader that selects only
declared partition identifiers and returns
\(X\in\mathbb{N}^{R\times K}\) plus aligned metadata.}{Assert shape, dtype,
non-negativity, category order, and exact selected dates. Requesting an
undeclared partition must fail.}

\gate{The smoke and configured captures produce schema-valid, time-indexed,
checksum-backed matrices with correct silence, missingness, and partition
semantics.}

\section{Stage 4: Dirichlet--Multinomial Model}

\task{P2.01}{Define estimator input contract}{Accept only two-dimensional
non-negative integer matrices with \(R>1\), \(K>1\), stable category order, and
a documented policy for zero-total rows and all-zero categories.}{Parameterized
tests reject every invalid dimension and policy violation; a valid matrix is
not mutated by training.}

\task{P2.02}{Verify the likelihood formula}{Implement the reference log
probability
\begin{align}
\log p(\mathbf{x}\mid\boldsymbol{\alpha})={}&
\log\Gamma(N+1)-\sum_{k=1}^{K}\log\Gamma(x_k+1)\\
&+\log\Gamma(\alpha_0)-\log\Gamma(N+\alpha_0)
+\sum_{k=1}^{K}\left[\log\Gamma(x_k+\alpha_k)-\log\Gamma(\alpha_k)\right],
\end{align}
where \(\alpha_0=\sum_k\alpha_k\). Make inclusion of the multinomial
coefficient explicit.}{Compare small cases with direct arbitrary-precision
probability enumeration and verify the coefficient convention with a
hand-calculated example.}

\task{P2.03}{Implement robust initialization}{Provide mean-concentration and
moment-based initializers, floors for sparse categories, and deterministic
fallback behavior.}{Test balanced, highly skewed, sparse, and nearly constant
synthetic matrices; every returned \(\alpha_k\) must be finite and strictly
positive.}

\task{P2.04}{Harden the fixed-point update}{Separate one Minka update from the
iteration loop, validate numerator and denominator, guard positivity, and
detect non-finite or zero updates.}{Unit-test the update against a stored
high-precision example and force each guard branch with synthetic pathological
inputs.}

\task{P2.05}{Define convergence criteria}{Require both parameter stability and
likelihood stability, use absolute plus relative tolerances, impose maximum
iterations, and distinguish converged, stalled, diverged, and exhausted
outcomes.}{Construct fixtures for each outcome and assert status, iteration
count, and diagnostic reason. Never label maximum-iteration exhaustion as
convergence.}

\task{P2.07}{Activate the modeling stage}{Load only the fit partition,
initialize and fit the estimator, derive \(\alpha_0\),
\(p_k=\alpha_k/\alpha_0\), and \(\psi=1/\alpha_0\), then persist diagnostics.}
{Run the CLI on the smoke matrix; assert a valid model artifact, successful
exit, correct categories, positive alpha, \(\sum_kp_k=1\), and no access to
calibration/test partitions.}

\task{P2.08}{Define the model artifact}{Include schema version, model version,
categories, parameters, training range, dataset/config/code hashes, backend,
convergence, runtime, and integrity checksum.}{Validate against schema,
round-trip exactly, reject category reordering and checksum corruption, and
verify no raw packet data is embedded.}

\task{P2.09}{Recover synthetic parameters}{Sample multiple seeded datasets from
known \(\boldsymbol{\alpha}^{*}\), fit them, and quantify error in mean
composition and concentration.}{Across increasing sample sizes, demonstrate
decreasing median estimation error or explain deviations; enforce predeclared
tolerances on a stable medium-size fixture.}

\gate{Partition-safe training produces an integrity-checked converged model and
recovers known synthetic parameters within declared tolerances.}

\section{Stage 5: LM Likelihood Backend}

\task{P3.01}{Define the backend protocol}{Create a stateless
\artifact{LikelihoodBackend} interface returning value, precision status, error
bound when available, backend version, and timing metadata; support
single-window and batch calls.}{Contract-test a fake backend and every real
backend for shape, dtype, metadata, immutability, error semantics, and
consistent coefficient convention.}

\task{P3.02}{Implement the standard gammaln backend}{Move SciPy-based
likelihood evaluation behind the interface and support stable vectorized batch
evaluation.}{Compare with the P2 arbitrary-precision oracle across small random
cases and with regression fixtures.}

\task{P3.03}{Build a high-precision oracle}{Use \artifact{mpmath} at
configurable precision to evaluate selected log probabilities and paired
log-gamma differences outside production code.}{Validate oracle precision by
recomputing at doubled digits; accepted reference digits must remain unchanged.
Store reference vectors and provenance.}

\task{P3.05}{Implement paired log-gamma differences}{Implement the LM primitive
for expressions such as
\(\log\Gamma(x+a)-\log\Gamma(a)\), with no silent fallback and explicit error
estimates.}{Compare a logarithmic grid of small/large \(x\), small/large \(a\),
and low-overdispersion cases with the doubled-precision oracle. Enforce
requested error bounds.}

\task{P3.06}{Assemble the Python LM backend}{Use the LM primitive for all
paired terms in the Dirichlet--Multinomial likelihood and retain an
exact/reference path for factorial terms.}{Contract-test against standard and
oracle backends over randomized and adversarial grids. Assert correct precision
status propagation.}

\task{P3.08}{Test extreme numerical regimes}{Cover zero counts, huge counts,
highly imbalanced categories, \(\alpha_k\) near the lower bound, very large
concentration, and cancellation-prone low-overdispersion cases.}{No accepted
result may be NaN or infinite. Compare to the oracle and enforce tolerances
appropriate to each regime.}

\task{P3.10}{Implement backend selection}{Select backend only through validated
configuration and record backend name/version in every model, score, and
manifest.}{Run identical scoring with each enabled backend and verify artifact
metadata. Invalid or unavailable selections must fail before data processing.}

\gate{The standard and LM backends satisfy one interface, agree with the
independent oracle within declared tolerances, and expose precision failures.}

\section{Stage 6: Calibration and Anomaly Output}

\task{P4.01}{Enforce fit/calibration isolation}{Load calibration windows
through a partition-scoped API that cannot return fit or final-test rows.}
{Assert disjoint identifiers and monkeypatch the fit loader to fail if
calibration invokes it. Store partition fingerprints in the threshold artifact.}

\task{P4.02}{Compute raw window scores}{Return one log probability
\(s_t=\log p(\mathbf{x}_t\mid\hat{\boldsymbol{\alpha}})\) and precision metadata
per eligible window.}{Compare each score with the backend single-window call;
preserve timestamps/order and explicitly classify missing windows.}

\task{P4.03}{Define normalized scores}{Implement a documented normalization,
such as per-packet score \(s_t/\max(N_t,1)\), while treating silent windows
according to policy rather than dividing silently.}{Test totals \(N_t=0,1\),
and large \(N_t\); assert finite values or explicit non-applicability and record
the formula in artifacts.}

\task{P4.05}{Calibrate empirical quantiles}{Compute the lower-tail threshold
\(\tau=Q_q(\{s_i\}_{i=1}^{m})\) with declared quantile convention and target
calibration false-positive rate.}{Compare with an independent quantile
implementation, test tiny/tied samples, and verify the observed calibration
alarm rate matches the convention within finite-sample limits.}

\task{P4.07}{Persist threshold artifacts}{Store model checksum, score type,
quantile, threshold, calibration range/fingerprint, sample count, backend, and
schema version.}{Round-trip and checksum-test the artifact. Scoring must reject
a threshold created for another model, backend convention, or category schema.}

\task{P4.08}{Implement online scoring}{Consume windows in timestamp order and
emit score records without future context. Keep state bounded and deterministic.}
{Replay the same stream in different chunk sizes and assert identical ordered
outputs; monitor memory over a long synthetic stream.}

\task{P4.09}{Implement anomaly decisions}{Apply the strict documented rule
\(s_t<\tau\), including explicit equality behavior, and retain both score and
threshold.}{Test values below, equal to, and above threshold, including
floating-point neighbors. Decisions must be backend-metadata aware.}

\task{P4.11}{Compute category residual explanations}{For \(N>0\), calculate
\[
r_k=\frac{x_k}{N}-\frac{\alpha_k}{\alpha_0},\qquad
e_k=N\frac{\alpha_k}{\alpha_0},
\]
and rank positive and negative deviations; use a separate explanation for
silence.}{Hand-check a known vector, assert residuals sum to zero within
tolerance, category ordering is deterministic under ties, and silent windows
never divide by zero.}

\task{P4.13}{Define anomaly event schemas}{Create a versioned event containing
device/window IDs, counts, score, threshold, decision, explanation, precision
status, model/threshold versions, and run ID.}{Schema-test mandatory fields,
reject incompatible category order, round-trip JSON/CSV, and verify no packet
payload or unnecessary endpoint identifier is present.}

\task{P0.12}{Create the end-to-end smoke command}{Provide one command that
validates configuration, processes the smoke fixture, trains, calibrates,
scores, and writes a report.}{Execute from a clean environment; assert
successful exit, expected artifact set, valid schemas and checksums,
deterministic decisions on repeated runs, and no modification of inputs.}

\gate{One command takes the smoke fixture through preprocessing, training,
LM-backed scoring, threshold calibration, and schema-valid anomaly events. This
is the working-prototype milestone.}

\section{Deferred Until After the Working Prototype}

The omitted roadmap tasks fall into the following later categories:
\begin{itemize}
  \item \textbf{Additional engineering assurance:} broader CI, coverage,
  property-based testing, structured logging, container hardening, governance,
  threat modeling, and expanded reports.
  \item \textbf{Research validation:} full anomaly-injection suite, detector
  baselines, uncertainty estimates, sensitivity/ablation studies, locked final
  evaluation, and publication figures.
  \item \textbf{Additional functionality:} adaptive retraining, forecasting,
  and advisory response support.
  \item \textbf{Deployment and polishing:} optimized native bindings, edge
  benchmarks, supply-chain controls, restart/exhaustion testing, deployment
  documentation, and release packaging.
\end{itemize}

After this document's final gate passes, return to the Atomic Implementation
Roadmap and complete the deferred tasks required by the thesis scope before
making research-quality, security, or deployment-readiness claims.

\end{document}
